\documentclass[12pt]{ximera}
\title{Homework 7: Combinations and Counting in Probability}
\begin{document}
\begin{abstract}
Sections 8.2--8.3: telling combinations from permutations, counting selections for
several gymnastics apparatuses at once, and turning counts into probabilities.
\end{abstract}
\maketitle

\section*{Sections 8.2--8.3: Combinations and Probability}

\begin{problem}
Identify whether each of the following can be counted using combinations or
permutations.

\begin{prompt}
\textbf{(a)} A track coach needs to select athletes for the first, second, third, and
fourth legs of a $4\times 100\text{m}$ relay from a team of $12$ runners.

\begin{multipleChoice}
\choice{Combinations}
\choice[correct]{Permutations}
\end{multipleChoice}

\begin{feedback}
Permutations: the four legs are different positions, so the order of selection
matters.
\end{feedback}
\end{prompt}

\begin{prompt}
\textbf{(b)} A PhD student has to select $5$ professors for their graduate committee
from among the $20$ professors in the department.

\begin{multipleChoice}
\choice[correct]{Combinations}
\choice{Permutations}
\end{multipleChoice}

\begin{feedback}
Combinations: a committee is just a group of five, with no roles to tell them apart.
\end{feedback}
\end{prompt}

\begin{prompt}
\textbf{(c)} Guessing the correct birth order of ten siblings.

\begin{multipleChoice}
\choice{Combinations}
\choice[correct]{Permutations}
\end{multipleChoice}

\begin{feedback}
Permutations: a birth order is an ordering of all ten, and different orders are
different guesses.
\end{feedback}
\end{prompt}

\begin{prompt}
\textbf{(d)} Playing a game of Memory, you choose which two cards you're going to flip
from a $3\times 4$ array of cards.

\begin{multipleChoice}
\choice[correct]{Combinations}
\choice{Permutations}
\end{multipleChoice}

\begin{feedback}
Combinations: what matters is \emph{which} pair of cards you turn over, not which one
you touched first. There are $C(12,2)=66$ possible pairs.
\end{feedback}
\end{prompt}
\end{problem}

\begin{problem}
In the 2024 Olympics, Team USA's roster listed $7$ women's gymnastics athletes. For the
team final, three gymnasts from each team compete on each apparatus. Women's gymnastics
has $4$ apparatuses, so for each one a coach selects $3$ gymnasts from the team.

\begin{prompt}
\textbf{(a)} For a single apparatus, how many ways were there to select three
athletes?

$\answer{35}$

\begin{feedback}
\[ C(7,3)=\frac{7\cdot 6\cdot 5}{3\cdot 2\cdot 1}=35. \]
\end{feedback}
\end{prompt}

\begin{prompt}
\textbf{(b)} How many ways were there to select athletes for all four apparatuses?

$\answer{1500625}$

\begin{hint}
A gymnast may compete on more than one apparatus, so the four selections are
independent.
\end{hint}

\begin{feedback}
Each apparatus independently gets one of the $35$ trios:
\[ 35^4 = 1{,}500{,}625. \]
\end{feedback}
\end{prompt}

\begin{prompt}
\textbf{(c)} In reality, only five of the athletes represented the team on any
apparatus. How many ways were there to select athletes for all four apparatuses from
only those five?

$\answer{10000}$

\begin{feedback}
Now each apparatus gets one of $C(5,3)=10$ trios, so there are $10^4=10{,}000$ ways.
\end{feedback}
\end{prompt}

\begin{prompt}
\textbf{(d)} If the selections were made at random from all seven athletes, what would
be the probability that every apparatus draws only from those five? Enter a fraction in
lowest terms.

$\answer{16/2401}$

\begin{feedback}
Divide part (c) by part (b):
\[
\frac{10{,}000}{1{,}500{,}625}=\left(\frac{10}{35}\right)^4=\left(\frac27\right)^4=\frac{16}{2401}\approx 0.0067,
\]
about $0.67\%$. Each apparatus independently has a $\frac27$ chance of avoiding the two
other gymnasts, and all four must.
\end{feedback}
\end{prompt}
\end{problem}

\begin{problem}
Determine the probability of each of the following, using counting techniques. Enter
fractions in lowest terms.

\begin{prompt}
\textbf{(a)} A magician draws $4$ cards from a shuffled deck of $52$. What is the
probability that they are the four aces?

The probability is $1$ in $\answer{270725}$.

\begin{feedback}
Exactly one of the $C(52,4)=270{,}725$ possible four-card hands is the four aces.
\end{feedback}
\end{prompt}

\begin{prompt}
\textbf{(b)} A photo is being taken of $5$ siblings of different heights, with two in
front and three behind. If the placement is chosen at random, what is the probability
that the shortest two are the ones in front?

$\answer{1/10}$

\begin{hint}
It only matters \emph{which} two siblings are in front.
\end{hint}

\begin{feedback}
There are $C(5,2)=10$ equally likely choices for the front pair, and exactly one of
them is the shortest two, so the probability is $\frac{1}{10}$.

(Counting full arrangements gives the same answer: $2!\cdot 3!=12$ of the $5!=120$
arrangements put the shortest two in front, and $\frac{12}{120}=\frac1{10}$.)
\end{feedback}
\end{prompt}

\begin{prompt}
\textbf{(c)} In a simplified BINGO game there are $25$ uniquely numbered balls, and your
card holds all $25$ numbers in a $5\times 5$ grid. What is the probability that the first
five balls drawn complete a row, a column, or a diagonal on your card?

$\answer{2/8855}$

\begin{hint}
How many rows, columns, and diagonals are there? And how many different sets of five
balls could be drawn first?
\end{hint}

\begin{feedback}
A winning set of five is one of the $5$ rows, $5$ columns, or $2$ diagonals --- $12$
sets in all. The first five balls form one of $C(25,5)=53{,}130$ equally likely sets.
So
\[ P=\frac{12}{53{,}130}=\frac{2}{8855}\approx 0.00023. \]
\end{feedback}
\end{prompt}
\end{problem}

\end{document}
